% TOP_PROBLEM: {"id": 3000046, "title": "Infinite Lucchesi-Younger", "category_id": 2, "category": {"id": 2, "name": "combinatorics", "display_name": "Combinatorics", "description": "Counting problems, graph theory, discrete structures.", "slug": "combinatorics", "order_index": 2, "created_at": "2026-07-31T15:26:25.671Z"}, "status": "open", "problem_number": "AMR-029-0046", "set_id": 13}
% FIRST_POSTED: 2026-10-02
% ENTRY_KIND: statement_only
% UnsolvedMath ID 3000046; source catalogue code AMR-029-0046
% !TeX program = lualatex
% Definition and sources only; this file is not an LLM solution attempt.
\documentclass[11pt,a4paper]{article}
\usepackage[margin=25mm]{geometry}
\usepackage{fontspec}
\setmainfont{lmroman10-regular.otf}[BoldFont=lmroman10-bold.otf,
 ItalicFont=lmroman10-italic.otf,BoldItalicFont=lmroman10-bolditalic.otf]
\usepackage{amsmath,amssymb,mathtools}
\usepackage{unicode-math}
\setmathfont{latinmodern-math.otf}
\AtBeginDocument{\let\setminus\smallsetminus}
\usepackage{xurl,hyperref}
\hypersetup{colorlinks=true,urlcolor=blue}
\setcounter{secnumdepth}{0}
\setlength{\parindent}{0pt}
\setlength{\parskip}{5pt}
\setlength{\emergencystretch}{3em}
\begin{document}
\section*{Infinite Lucchesi-Younger}\label{3000046}
{\small UnsolvedMath ID 3000046 · AMR-029-0046 · Combinatorics · AMR Open Problem Lists}
\subsection{Definitions and mathematical statement}
Let $D=(V,A)$ be a digraph, possibly infinite and with parallel arcs
distinguished. For $X\subseteq V$, write $\delta^+(X)$ and
$\delta^-(X)$ for the arcs leaving and entering $X$, respectively.
A dicut is a nonempty set $\delta^+(X)$ for which
$\delta^-(X)=\varnothing$. A finite dicut has finitely many arcs;
its defining vertex set need not be finite.

Does every such digraph admit a family $\mathcal C$ of pairwise
arc-disjoint finite dicuts and a choice $a_C\in C$ for each
$C\in\mathcal C$ such that
\[
                  F=\{a_C:C\in\mathcal C\}
\]
meets every finite dicut of $D$? Pairwise disjointness ensures that
the chosen representatives are distinct and $|F\cap C|=1$ for each
member of the packing. No additional arcs are allowed in $F$.

The hitting requirement ranges over all finite dicuts, including those
not in $\mathcal C$. The family and its representative set may both
be infinite. If there are no finite dicuts, the empty family and empty
set satisfy the requirement.

For finite digraphs this is the packing/transversal conclusion of the
Lucchesi--Younger theorem. In the infinite setting an equality of
cardinalities alone is weaker than the requested compatible packing
and transversal. Neither local finiteness nor countability is assumed,
and infinite dicuts are deliberately excluded from the family to be hit.
\subsection{Short English statement}
Can finite directed cuts in every infinite digraph be packed with one representative from each packed cut forming a transversal of all finite directed cuts?
\subsection{Sources}
\begin{itemize}
\item[S1] ULAM.AI and the UnsolvedMath Contributors, supplied problems.json, record 3000046 (AMR-029-0046), AMR Open Problem Lists. Catalogue identity and source transcription; CC BY 4.0.
\url{https://huggingface.co/datasets/ulamai/UnsolvedMath}.
\item[S2] Egres Open - List of open problems, Open-problem page: Infinite Lucchesi-Younger. Original source indexed by AMR-029-0046.
\url{https://oldlemon.cs.elte.hu/egres/open/List_of_open_problems}.
\item[S3] EGRES Open, Infinite Lucchesi-Younger. Individual source statement and remarks.
\url{https://lemon.cs.elte.hu/egres/open/Infinite_Lucchesi-Younger}.
\end{itemize}
\end{document}
